\section{Conclusion}
\label{sec:conclusion}

Mixture-of-Models moves conditional computation from parameters within one
checkpoint to execution across independently versioned models. We define the
served object as a preference-conditioned distribution over finite, typed
traces, with hard admissibility, portable logical identity, and
environment-specific binding. Selection, cascades, fusion, multi-round
collaboration, and bounded workflows are model topologies rather than
application glue. The resulting intelligence is system-level: it arises from
co-designing the trace distribution with its physical realization, not from any
constituent or control component in isolation.

\vsr{} shows that an execution-facing subset of this abstraction can be
declarative and executable without joint access to constituent weights. The
portable bundle, lock, and conformance suite define its next implementation
boundary. The central empirical question is whether added model diversity
increases utility under matched active compute. If it does, model runtimes will
need to execute and evolve composite models across providers, jurisdictions,
devices, and fleets as reliably as they execute checkpoints across machines.
The resulting mission is to make virtual composite models buildable,
evaluable, portable, and servable with the same rigor as checkpoints.
